% !PRINTED 217
\weberpage{217}

Eine Invariante vom Grade $n(n-1)$ ist auch die Discriminante $\Delta$ der Gleichung (3), oder das Quadrat des Differenzenproductes

\weq{4}{\begin{gathered} \Pi = (y_1 - y_2)(y_1 - y_3)\cdots \\ (y_2-y_3)\cdots \end{gathered}}

Um über die Bildung dieser Grösse näheren Aufschluss zu bekommen, erinnern wir uns, dass wir $y$ erhalten, wenn wir in

\[ \frac{f(t)}{t-x} - \frac{1}{n}f'(t) \]

$t^k$ durch $t_k$ ersetzen. Wir erhalten also $y_1-y_2$ aus der Formel

\weq{5}{y_1-y_2 = (x_1-x_2)\frac{f(t)}{(t-x_1)(t-x_2)},}

wenn wir dieselbe Vertauschung machen.

Der Quotient

\[ \frac{f(t)}{(t-x_i)(t-x_k)} \]

ist eine ganze rationale Function von $t$ vom Grade $n-2$. Ersetzen wir darin $t^k$ durch $t_k$, so möge er in $t_{i,k}$ übergehen; wir haben dann

\weq{6}{y_i - y_k = (x_i-x_k)t_{i,k}.}

Wenn wir also

\weq{7}{\begin{gathered} a_0^{n-1}\Theta = t_{1,2}t_{1,3}\ldots t_{1,n} \\ t_{2,3}\ldots t_{2,n} \\ t_{n-1,n} \end{gathered}}

setzen, so ist $\Theta$ eine homogene ganze Function vom Grade $\frac{1}{2}n(n-1)$ in Bezug auf die $t$, und sie ist ausserdem als symmetrische Function der $x$ rational durch die Coëfficienten $a$ darstellbar.

Aus (6) folgt aber, wenn $D$, wie im \S.\,46, (3) die Discriminante von $f(x)$ bedeutet,

\weq{8}{\Delta = D\Theta^2,}

woraus zu schliessen ist, dass $\Theta$ eine Invariante ist, \gesperrt{die den Nenner $a_0$ nicht mehr enthält}. $\Delta$ ist in den Coëfficienten $a$ vom Grade $n(n-1)$, während $D$ nach \S.\,46 vom Grade $2n-2$ ist; daraus folgt, dass $\Theta$ vom Grade $\frac{1}{2}(n-1)(n-2)$ in den Coëfficienten $a$ ist.
